\section{Scope, Method, and Qualifications}

\subsection{Reader, question, and method}

This is a discursive theological article for a serious general reader, able to follow careful distinctions but not presumed to know scholastic terminology. It is Part II of a three-part series, \emph{The Creature Before God}; each part is a complete essay. Part I (\emph{Ontological Vertigo}) established received being, present conservation, the one-sided real relation, and the fork of recoil or consent; this essay inherits those conclusions without re-arguing them. Its governing question is what the consenting creature owes: in what sense a return is due to God when no equal return is possible, and what form the rendering takes in an embodied life. Its claim: the creature's situation before God has the form of debt without possibility of settlement; the tradition locates the rendering virtue---religion---among the moral virtues annexed to justice, falling short of justice's equality but not of its due; what is owed is acknowledgment rather than equivalence (equality of wills, not of things); religion's elicited acts (devotion, prayer, adoration, sacrifice) engage the body because man signifies through sensibles; the rendering profits the renderer alone, God receiving nothing; sacrifice's natural-law core is acknowledgment in sign, while revelation discloses the incarnate Son's unique saving oblation and creatures' dependent participation in it; the due return must be free because what is owed is consent itself; and no aggregate of fulfilled acts closes the account, the debt of gratitude belonging finally to charity---the threshold of Part III.

The argument proceeds in epistemic order: the psalmist's question and the two false exits; justice's definition and its failed equality toward God, with religion and gratitude analyzed from Aquinas; the ingratitude diagnosis of Romans 1 read textually; the virtue of religion's placement and its two ranges of acts; the elicited acts with their anti-magical boundary (ST II-II, q.~81, a.~7); sacrifice argued from natural reason and then Augustine's \emph{City of God} X.5--6, with the seam between definition and Christian confession expressly marked; the Church's teaching (CCC 2095--2109, \emph{Dei Filius} ch.~3, \emph{Dignitatis humanae} 1--3) cited as received teaching; and the closing synthesis on the open account. Revealed claims (Christ's oblation, the Church's self-offering, the sacrament of the altar, the obedience of faith) are labeled as received where they appear and are nowhere presented as conclusions of the philosophical thread. English is the working language; no liturgical rite, civil or canonical jurisdiction, or mutable discipline governs any conclusion; the canonical and liturgical determinations of worship (which rites, days, and obligations) are expressly outside scope. All online witnesses were checked on 2026-07-24.

\begin{threestable}{Source class}{Function in the article}{Governing boundary}
Scripture & Psalm 114--115 (Vulgate numbering) supplies the governing question and its answer; Romans 1:21, 25 the ingratitude diagnosis; Romans 12:1, 5 the living-sacrifice hinge inside Augustine's argument; James 1:27 cited through Aquinas. & Quoted from the public-domain Douay--Rheims (Challoner) at exact loci; readings are textual and modest; no independent exegetical apparatus is claimed.\\
Thomistic framework & ST II-II, qq.~58, 80--85, 106 govern the analysis of justice, religion, devotion, prayer, adoration, sacrifice, and gratitude; ST I, qq.~13, 20, 104 and II-II, qq.~84, 129, 161 are inherited from Part I. & Adopted as the article's framework where Catholic theology permits schools; not asserted as dogma; each locus used at its exact argumentative level.\\
Patristic witness & Augustine, \emph{De civitate Dei} X.5--6, supplies the sign-theology of sacrifice and the fountain/light images; X.3 and X.19 appear inside Aquinas's citations. & The X.6 climax (the redeemed city offered through the High Priest; the sacrament of the altar) is marked as Christian confession, not ascent.\\
Authoritative Catholic teaching & CCC 2095--2109 on the virtue of religion; \emph{Dei Filius} ch.~3 on dependence grounding the obedience of faith; \emph{Dignitatis humanae} 1--3 on the untouched duty and civil immunity from coercion. & Cited as the Church's teaching at its stated authority; \emph{Dignitatis humanae} is read for its stated ground (the free nature of religious acts) with juridical applications excluded.\\
Project synthesis & The reading of \latin{quid retribuam}/\latin{accipiam}; the ledger, pronoun, basket, window/wound, and counterfeit-currency images; ``recoil's liturgical form''; the ``spiritual bookkeeping'' diagnosis; all final formulations. & Original editorial synthesis argued from the checked sources; attributed to no source and creating no doctrine.\\
\end{threestable}

\subsection{Included and excluded scope}

Included are: the psalmist's question (Ps 115:12 Vulg.) with its answer in reception, invocation, vows, and public praise; justice as rendering the due by a constant and perpetual will (II-II.58.1); the annexation analysis (II-II.80.1) with its citation of the psalm; the parental analogue and the ranked causes of debt (II-II.106.1); the exceeding, limitless debt of gratitude (II-II.106.6 with ad~2); Romans 1:21, 25 as ingratitude preceding idolatry; religion's etymologies and definition (II-II.81.1), its status as virtue (81.2), special virtue (81.4), moral not theological virtue (81.5 with ad~1), chief of moral virtues (81.6); elicited versus commanded acts (81.1 ad~1) with James 1:27; devotion (82.1); petitionary prayer and providence (83.2); twofold adoration and the profession of prostration (84.2 with ad~2); worship's profit to the worshipper alone (81.7); sacrifice's natural-law core and positive determination (85.1 with ad~1); the outward sign of the soul's self-offering (85.2); Augustine on divine non-need, the visible sacrament of the invisible sacrifice, true sacrifice, mercy, and the redeemed city (\emph{De civ. Dei} X.5--6); CCC 2095--2109; \emph{Dei Filius} ch.~3; \emph{Dignitatis humanae} 1--3; and the closing analysis of the open account with its arrow toward charity.

Excluded are: any treatise-scale account of virtue, law, grace, merit, sacraments, liturgy, or ecclesiology; the history of sacrificial practice and of Sabbath or Sunday observance; every canonical determination of worship, including the Sunday obligation, feast days, and their excuses (no current-law claim is made anywhere); Eastern Catholic law and particular law; the Mass considered sacramentally and the theology of Eucharistic sacrifice beyond Augustine's and the Catechism's quoted sentences; vows and states of life beyond the Catechism's definition; the juridical interpretation and application history of \emph{Dignitatis humanae}; theodicy; merit and the supernatural quality of acts; interreligious questions and the status of non-Christian sacrifice beyond Aquinas's natural-law claim; ascetical and penitential practice; scrupulosity and pastoral counsel; and judgments about any person's obligations, culpability, or interior state.

\subsection{Material qualifications}

\begin{enumerate}[leftmargin=1.45em,itemsep=0.25em]
\item \emph{Ontological vertigo}, \emph{recoil}, and \emph{consent} are Part I's coinages, used here in the senses defined there; \emph{religion} is used throughout in the tradition's sense of a virtue of rendering, not the modern sense of a belief system, and the essay says so where the word is introduced.
\item ``Debt,'' ``ledger,'' ``account,'' and cognates are analogical throughout. The essay expressly denies that the creature's due is a closable obligation, that creation was a contract, or that God gains by the rendering; the analogical register is stated in sections 1--2 and governs every later use.
\item The claim that religion is a moral virtue annexed to justice, not a theological virtue, is Aquinas's (II-II.81.5) and is presented as the adopted framework of broad Thomistic tradition, not as dogma; the Catechism's placement of religion under the first commandment (CCC 2095--2096) is cited as its catechetical register.
\item The reading of Romans 1:21--25 asserts only Paul's stated order (failed glorification and thanksgiving preceding darkened thought and idolatry) and builds no claim about the salvation, culpability, or interior state of any historical people or person.
\item ``Recoil's liturgical form'' (idolatry as preference for a repayable god) is project synthesis offered as an interpretation of the Pauline sequence, not as exegetical consensus or a historical thesis about ancient religion.
\item Aquinas's natural-law claim about sacrifice (II-II.85.1) is generic: that acknowledgment in sensible sign is naturally fitting. The determination of rites is positive (85.1 ad~1); the essay makes no claim that any particular rite, victim, or observance is naturally required, and no comparative-religion thesis.
\item The efficacy denied of worship (adding nothing to God) and the efficacy affirmed (the perfection of the worshipper) are both Aquinas's (II-II.81.7); the essay's ``enacted truth'' formulation is its own gloss and claims nothing about sacramental efficacy \latin{ex opere operato}, which is outside scope.
\item Prayer is treated only at the joint of providence and secondary causality (II-II.83.2); the essay makes no claim about what may be asked, the certainty of being heard, or contemplative prayer, all reserved to their own treatments.
\item Augustine's X.6 climax and the Catechism's sentences on Christ's one perfect sacrifice are received Christian confession, labeled as such; the essay's philosophical thread establishes that acknowledgment in sign is owed, and neither reaches nor argues the Priest, the Passion, or the altar. The sacrament of the altar is named, not developed; Eucharistic theology belongs to Part III's territory and beyond.
\item \emph{Dei Filius} ch.~3 concerns the obedience of faith owed to God revealing; it is cited for its dogmatically stated premise (total dependence grounding total obligation), and the essay does not conflate the obedience of faith with the virtue of religion.
\item \emph{Dignitatis humanae} is cited at its stated level: a conciliar declaration on civil immunity from coercion that itself declares the traditional moral duty untouched. The essay takes no position on its juridical applications, on church--state arrangements, or on the declaration's reception debates, and does not derive the immunity from the moral analysis alone---the Council's own grounds are quoted.
\item The freedom required of religious acts concerns their constitution as acts of the virtue (a coerced act renders nothing); nothing here denies that obligations can bind in conscience, that parents rightly form children in worship, or that the Church may order her own common life---all outside scope.
\item ``Spiritual bookkeeping''---fulfilled acts construed as installments toward a discharged standing before God---is project synthesis naming a temptation; it neither denies the objective goodness of fulfilled religious acts nor addresses the theology of merit, satisfaction, or indulgences, which are excluded.
\item The closing claims about what the giver ``wants,'' about mutuality, desire, and union, are named as the matter of charity, revelation, and Part III; they are pointers beyond this essay's argument, not conclusions of it.
\item No pastoral, canonical, or liturgical advice is given, and no judgment is made about any person's obligation or culpability in worship.
\end{enumerate}

\subsection{Notes, translations, rights, currentness, and review}

The numbered Notes carry exact citations, edition and translation identifications, verification dates, and qualifications that would interrupt the argument; no indispensable premise lives only in a note. Scripture is quoted from the public-domain Douay--Rheims translation (Challoner revision), with the Vulgate Latin of Psalm 115:12--13 checked against a public Clementine text. Aquinas is quoted from the public-domain English Dominican translation with Latin checked at Corpus Thomisticum. Augustine is quoted from the public-domain Dods translation (1871), checked in the source library's tracked transcription of the 1871 printing at the registered passage for \emph{De civitate Dei} X.5--6 and against the New Advent transcription of the same translation. \emph{Dei Filius} is quoted from the Holy See's Latin web text with a labeled working gloss; \emph{Dignitatis humanae} and the Catechism are quoted briefly from the official Vatican English web texts with attribution. Official texts and third-party translations remain outside the project's CC~BY~4.0 grant; project-created prose, images, and organization are project content.

All online witnesses were checked on 2026-07-24; no mutable law or discipline is relied on, and no canonical obligation is asserted or interpreted. This revision received internal argumentative, source-consistency, quotation, rights, and production review by the authoring agent. Independent scriptural, patristic, Thomistic, theological, liturgical, literary, and ecclesiastical review is outstanding; no imprimatur, nihil obstat, or ecclesiastical approval is claimed, and internal checking is not independent review.
